% GENERATED FILE. Edit canonical lessons, then run npm run generate:books.
% canonical-source-hash: fnv1a64:c760672659fa41c4
% canonical-lessons: PA-C284-utarna, PA-C284-utaarna, PA-C284-dhakka, PA-C284-akhri, PA-C284-apparna

\chapter{Getting Off}
\label{ch:pa-a2-getting-off}
\hlchaptermodality{284}

\begin{chapteropening}
\textbf{By the end of this chapter:} \emph{I can say where I will get off.}

Complete the last lesson of chapter 284: I can say where I will get off.
\end{chapteropening}

\section[\texorpdfstring{\pa{ਉਤਰਨਾ} (utarnā)}{ਉਤਰਨਾ (utarnā)}]{\pa{ਉਤਰਨਾ} (utarnā) — to get off; to come down}

\label{lesson:PA-C284-utarna}



\begin{quote}
Before the new one: say the Punjabi for to get stuck, then the Punjabi for to sway.
\end{quote}

\subsection*{You'll want to know: \pa{ਉਤਰਨਾ}}
\textbf{\pa{ਉਤਰਨਾ}} — \emph{utarnā} — \textquotedblleft{}to get off; to come down\textquotedblright{}.

\begin{quote}
\pa{ਮੈਂ} \pa{ਅਗਲੇ} \pa{ਸਟੇਸ਼ਨ} \pa{ਉੱਤੇ} \pa{ਉਤਰਾਂਗਾ।}
\end{quote}

\begin{quote}
\emph{maiṁ agle saṭeshan utte utrāṅgā.}
\end{quote}

\begin{quote}
I will get off at the next station.
\end{quote}

Past (I …): \pa{ਮੈਂ} \pa{ਬੱਸ} \pa{ਤੋਂ} \pa{ਉਤਰਿਆ}, \emph{maiṁ bass ton utriā}.

\begin{cousinweb}
\textbf{\pa{ਉਤਰਨਾ}} (\emph{utarnā}) is inherited from Sanskrit \emph{uttarati}, \textquotedblleft{}alights, descends\textquotedblright{}; CDIAL lists Punjabi \emph{uttarṇā}, \textquotedblleft{}to descend\textquotedblright{}, under it (CDIAL 1770).
\end{cousinweb}

\begin{grammarlens}[title={What you've built}]
One more word for saying what comes next.
\end{grammarlens}

\subsection*{Your turn}
\begin{itemize}
\raggedright
  \item \emph{Say it:} \emph{utarnā}
  \item \emph{Say it:} \emph{utarnā}, once more
  \item \emph{Say it:} \emph{maiṁ agle saṭeshan utte utrāṅgā.} — I will get off at the next station.
  \item \emph{Recall:} say the Punjabi for to get stuck, then the Punjabi for to sway, then say \emph{utarnā} again
\end{itemize}

\begin{tcolorbox}[breakable,colback=teal!4,colframe=teal!35!black,title={Before you move on}]
What does \pa{ਉਤਰਨਾ} mean? (\textquotedblleft{}To get off; to come down\textquotedblright{}.) Say it once more.
\end{tcolorbox}
\section[\texorpdfstring{\pa{ਉਤਾਰਨਾ} (utārnā)}{ਉਤਾਰਨਾ (utārnā)}]{\pa{ਉਤਾਰਨਾ} (utārnā) — to take down; to unload}

\label{lesson:PA-C284-utaarna}



\begin{quote}
Before the new one: say the Punjabi for to sway, then the Punjabi for to get off.
\end{quote}

\subsection*{You'll want to know: \pa{ਉਤਾਰਨਾ}}
\textbf{\pa{ਉਤਾਰਨਾ}} — \emph{utārnā} — \textquotedblleft{}to take down; to unload\textquotedblright{}.

\begin{quote}
\pa{ਅਮਨ} \pa{ਬੱਸ} \pa{ਦੀ} \pa{ਛੱਤ} \pa{ਤੋਂ} \pa{ਅਟੈਚੀ} \pa{ਉਤਾਰੇਗਾ।}
\end{quote}

\begin{quote}
\emph{Aman bass dī chhatt ton aṭaichī utāregā.}
\end{quote}

\begin{quote}
Aman will take the suitcase down from the roof of the bus.
\end{quote}

Past (I …): \pa{ਮੈਂ} \pa{ਅਟੈਚੀ} \pa{ਉਤਾਰੀ}, \emph{maiṁ aṭaichī utārī}.

\begin{cousinweb}
\textbf{\pa{ਉਤਾਰਨਾ}} (\emph{utārnā}) is the causative partner of \emph{utarnā}, inherited from Sanskrit \emph{uttārayati}, \textquotedblleft{}makes descend\textquotedblright{}; CDIAL lists Punjabi \emph{utārṇā}, \textquotedblleft{}to make descend\textquotedblright{}, under it (CDIAL 1770).
\end{cousinweb}

\begin{grammarlens}[title={What you've built}]
One more word for saying what comes next.
\end{grammarlens}

\subsection*{Your turn}
\begin{itemize}
\raggedright
  \item \emph{Say it:} \emph{utārnā}
  \item \emph{Say it:} \emph{utārnā}, once more
  \item \emph{Say it:} \emph{Aman bass dī chhatt ton aṭaichī utāregā.} — Aman will take the suitcase down from the roof of the bus.
  \item \emph{Recall:} say the Punjabi for to sway, then the Punjabi for to get off, then say \emph{utārnā} again
  \item \emph{Write it:} copy \textbf{\pa{ਉਤਾਰਨਾ}} once, checking the whole word
\end{itemize}

\begin{tcolorbox}[breakable,colback=teal!4,colframe=teal!35!black,title={Before you move on}]
What does \pa{ਉਤਾਰਨਾ} mean? (\textquotedblleft{}To take down; to unload\textquotedblright{}.) Say it once more.
\end{tcolorbox}
\section[\texorpdfstring{\pa{ਧੱਕਾ} (dhakkā)}{ਧੱਕਾ (dhakkā)}]{\pa{ਧੱਕਾ} (dhakkā) — a push; a shove}

\label{lesson:PA-C284-dhakka}



\begin{quote}
Before the new one: say the Punjabi for to get off, then the Punjabi for to take down.
\end{quote}

\subsection*{You'll want to know: \pa{ਧੱਕਾ}}
\textbf{\pa{ਧੱਕਾ}} — \emph{dhakkā} — \textquotedblleft{}a push; a shove\textquotedblright{}.

\begin{quote}
\pa{ਅਮਨ} \pa{ਨੇ} \pa{ਦਰਵਾਜ਼ੇ} \pa{ਨੂੰ} \pa{ਧੱਕਾ} \pa{ਦਿੱਤਾ।}
\end{quote}

\begin{quote}
\emph{Aman ne darvāze nū̃ dhakkā dittā.}
\end{quote}

\begin{quote}
Aman gave the door a push.
\end{quote}

A masculine noun.

\begin{cousinweb}
\textbf{\pa{ਧੱਕਾ}} (\emph{dhakkā}) is inherited from a reconstructed \emph{*dhakk}, \textquotedblleft{}push, strike\textquotedblright{}; CDIAL lists Punjabi \emph{dhakkā}, \textquotedblleft{}a push, blow\textquotedblright{}, under it (CDIAL 6701).
\end{cousinweb}

\begin{grammarlens}[title={What you've built}]
One more word for saying what comes next.
\end{grammarlens}

\subsection*{Your turn}
\begin{itemize}
\raggedright
  \item \emph{Say it:} \emph{dhakkā}
  \item \emph{Say it:} \emph{dhakkā}, once more
  \item \emph{Say it:} \emph{Aman ne darvāze nū̃ dhakkā dittā.} — Aman gave the door a push.
  \item \emph{Recall:} say the Punjabi for to get off, then the Punjabi for to take down, then say \emph{dhakkā} again
  \item \emph{Write it:} copy \textbf{\pa{ਧੱਕਾ}} once, checking the whole word
\end{itemize}

\begin{tcolorbox}[breakable,colback=teal!4,colframe=teal!35!black,title={Before you move on}]
What does \pa{ਧੱਕਾ} mean? (\textquotedblleft{}A push; a shove\textquotedblright{}.) Say it once more.
\end{tcolorbox}
\section[\texorpdfstring{\pa{ਆਖ਼ਰੀ} (ākhrī)}{ਆਖ਼ਰੀ (ākhrī)}]{\pa{ਆਖ਼ਰੀ} (ākhrī) — last; final}

\label{lesson:PA-C284-akhri}



\begin{quote}
Before the new one: say the Punjabi for to take down, then the Punjabi for a push.
\end{quote}

\subsection*{You'll want to know: \pa{ਆਖ਼ਰੀ}}
\textbf{\pa{ਆਖ਼ਰੀ}} — \emph{ākhrī} — \textquotedblleft{}last; final\textquotedblright{}.

\begin{quote}
\pa{ਇਹ} \pa{ਆਖ਼ਰੀ} \pa{ਸਟੇਸ਼ਨ} \pa{ਹੈ।}
\end{quote}

\begin{quote}
\emph{eh ākhrī saṭeshan hai.}
\end{quote}

\begin{quote}
This is the last station.
\end{quote}

\begin{cousinweb}
\textbf{\pa{ਆਖ਼ਰੀ}} (\emph{ākhrī}) is a loan, usually traced through Persian to Arabic \emph{ākhir}, \textquotedblleft{}last\textquotedblright{}, with the ending \emph{-ī}.
\end{cousinweb}

\begin{grammarlens}[title={What you've built}]
One more word for saying what comes next.
\end{grammarlens}

\subsection*{Your turn}
\begin{itemize}
\raggedright
  \item \emph{Say it:} \emph{ākhrī}
  \item \emph{Say it:} \emph{ākhrī}, once more
  \item \emph{Say it:} \emph{eh ākhrī saṭeshan hai.} — This is the last station.
  \item \emph{Recall:} say the Punjabi for to take down, then the Punjabi for a push, then say \emph{ākhrī} again
  \item \emph{Write it:} copy \textbf{\pa{ਆਖ਼ਰੀ}} once, checking the whole word
\end{itemize}

\begin{tcolorbox}[breakable,colback=teal!4,colframe=teal!35!black,title={Before you move on}]
What does \pa{ਆਖ਼ਰੀ} mean? (\textquotedblleft{}Last; final\textquotedblright{}.) Say it once more.
\end{tcolorbox}
\section[\texorpdfstring{\pa{ਅੱਪੜਨਾ} (appaṛnā)}{ਅੱਪੜਨਾ (appaṛnā)}]{\pa{ਅੱਪੜਨਾ} (appaṛnā) — to get to; to reach a place}

\label{lesson:PA-C284-apparna}



\begin{quote}
Before the new one: say the Punjabi for a push, then the Punjabi for last.
\end{quote}

\subsection*{You'll want to know: \pa{ਅੱਪੜਨਾ}}
\textbf{\pa{ਅੱਪੜਨਾ}} — \emph{appaṛnā} — \textquotedblleft{}to get to; to reach a place\textquotedblright{}.

\begin{quote}
\pa{ਬੱਸ} \pa{ਦੁਪਹਿਰ} \pa{ਤੱਕ} \pa{ਸ਼ਹਿਰ} \pa{ਅੱਪੜੇਗੀ।}
\end{quote}

\begin{quote}
\emph{bass dupahir tak shahir appṛegī.}
\end{quote}

\begin{quote}
The bus will get to the city by noon.
\end{quote}

Past (I …): \pa{ਮੈਂ} \pa{ਸਟੇਸ਼ਨ} \pa{ਅੱਪੜਿਆ}, \emph{maiṁ saṭeshan appṛiā}.

\begin{cousinweb}
\textbf{\pa{ਅੱਪੜਨਾ}} (\emph{appaṛnā}) is inherited from Sanskrit \emph{āpatati}, \textquotedblleft{}flies towards\textquotedblright{}; CDIAL lists Punjabi \emph{appaṛṇā}, \textquotedblleft{}to arrive, come to hand\textquotedblright{}, under it (CDIAL 1193).
\end{cousinweb}

\begin{grammarlens}[title={What you've built}]
One more word for saying what comes next.
\end{grammarlens}

\subsection*{Your turn}
\begin{itemize}
\raggedright
  \item \emph{Say it:} \emph{appaṛnā}
  \item \emph{Say it:} \emph{appaṛnā}, once more
  \item \emph{Say it:} \emph{bass dupahir tak shahir appṛegī.} — The bus will get to the city by noon.
  \item \emph{Recall:} say the Punjabi for a push, then the Punjabi for last, then say \emph{appaṛnā} again
  \item \emph{Write it:} copy \textbf{\pa{ਅੱਪੜਨਾ}} once, checking the whole word
\end{itemize}

\begin{tcolorbox}[breakable,colback=teal!4,colframe=teal!35!black,title={Before you move on}]
What does \pa{ਅੱਪੜਨਾ} mean? (\textquotedblleft{}To get to; to reach a place\textquotedblright{}.) Say it once more.
\end{tcolorbox}
